% Phasor SSM Derivation — technical derivation for inclusion in a larger paper
% Compile standalone: pdflatex phasor_ssm_derivation.tex

\documentclass[11pt]{article}
\usepackage{amsmath,amssymb,amsthm}
\usepackage{bm}
\usepackage{booktabs}
\usepackage{enumitem}
\usepackage[margin=1in]{geometry}
\usepackage{hyperref}

\newcommand{\R}{\mathbb{R}}
\newcommand{\C}{\mathbb{C}}
\newcommand{\im}{\mathrm{i}}
\newcommand{\dd}{\mathrm{d}}
\newcommand{\dt}{\Delta t}
\newcommand{\eg}{e.g.}
\newcommand{\ie}{i.e.}

\theoremstyle{definition}
\newtheorem{definition}{Definition}
\newtheorem{proposition}{Proposition}
\newtheorem{remark}{Remark}

\title{The Phasor State-Space Model:\\
A Derivation from Resonate-and-Fire Dynamics\\
to Holographic Distributed Representations}
\author{}
\date{}

\begin{document}
\maketitle

%======================================================================
\section{The Resonate-and-Fire Neuron}\label{sec:rf}
%======================================================================

\subsection{Single-neuron dynamics}

A Resonate-and-Fire (R\&F) neuron is a damped complex oscillator driven by
input current~$I(t)$:
\begin{equation}\label{eq:rf-ode}
  \frac{\dd z}{\dd t} = k\,z(t) + w\,I(t),
  \qquad k = \lambda + \im\omega,
\end{equation}
where $z(t)\in\C$ is the complex membrane potential,
$\lambda<0$ is the leakage (decay rate),
$\omega = 2\pi/T$ is the angular frequency with period~$T$, and
$w\in\R$ is a synaptic weight.
This is a one-dimensional complex linear ODE---equivalently, a
two-dimensional real linear system.

The potential decomposes into amplitude and phase:
\begin{equation}
  z(t) = r(t)\,e^{\im\pi\theta(t)},
  \qquad r(t) = |z(t)|,
  \quad \theta(t) = \arg(z(t))/\pi \in [-1,1].
\end{equation}
The neuron fires when $\operatorname{Im}(z)$ reaches a local maximum,
\ie, when the phase aligns with the natural oscillation.
Between spikes, the amplitude decays exponentially ($r\propto e^{\lambda t}$),
so \textbf{phase is the information-carrying variable}---amplitude merely
tracks signal strength.

\subsection{Phase as the natural observable}

In a biologically motivated network, the internal potential~$z$ of a neuron
is not directly observable by other neurons. Other neurons receive only
\textbf{spikes}---discrete timing events that encode the phase of~$z$ at
peak voltage. The natural output is therefore
\begin{equation}\label{eq:phase-readout}
  \theta = \arg(z)/\pi \in [-1,1],
\end{equation}
which can be extracted from the complex potential via normalization to the
unit circle:
\[
  \theta = \arg\!\bigl(z/|z|\bigr)/\pi.
\]
This discards amplitude and isolates the quantity that downstream neurons
can reconstruct from spike timing.

\subsection{Multi-channel R\&F layer}\label{sec:rf-multichannel}

A layer of $C$ R\&F neurons, each receiving a weighted projection of
$C_{\mathrm{in}}$ input channels, evolves as
\begin{equation}\label{eq:rf-layer}
  \frac{\dd z_c}{\dd t}
    = k_c\,z_c(t) + \sum_{j=1}^{C_{\mathrm{in}}} W_{cj}\,I_j(t),
  \qquad c=1,\dots,C,
\end{equation}
where $k_c = \lambda_c + \im\omega_c$ is a per-channel complex eigenvalue
and $W\in\R^{C\times C_{\mathrm{in}}}$ is the weight matrix.
Each neuron is an independent damped oscillator; cross-channel coupling
occurs only through~$W$.

This diagonal structure is the \textbf{physically natural} architecture:
each oscillator is a resonator tuned to its own $(\lambda_c,\omega_c)$.
The collection of oscillators forms a filter bank analogous to cochlear
hair cells or cortical oscillatory populations (theta, alpha, gamma bands).

%======================================================================
\section{State-Space Models}\label{sec:ssm}
%======================================================================

\subsection{The continuous SSM}

A continuous-time linear state-space model is defined by
\begin{equation}\label{eq:ssm-continuous}
  \frac{\dd u}{\dd t} = A\,u(t) + B\,x(t),
  \qquad y(t) = C\,u(t) + D\,x(t),
\end{equation}
where $u(t)\in\R^N$ (or $\C^N$) is the hidden state, $x(t)$ is the
input, $y(t)$ is the output, and $A,B,C,D$ are parameter matrices of
appropriate dimension.

\subsection{Three computational views}

The key insight of modern SSMs (S4, S4D, Mamba) is that this system admits
three equivalent computational forms:
\begin{enumerate}[nosep]
  \item \textbf{Recurrent.}
    Via zero-order-hold (ZOH) discretization with step~$\dt$:
    \[
      u[n+1] = \bar A\,u[n] + \bar B\,x[n],
      \qquad y[n] = C\,u[n] + D\,x[n],
    \]
    where $\bar A = e^{A\dt}$ and $\bar B = (e^{A\dt}-I)A^{-1}B$.

  \item \textbf{Convolutional.}
    Unrolling the recurrence yields a causal convolution:
    $y = K * x$, where $K[n] = C\,\bar A^n\,\bar B$ is the impulse-response
    kernel.

  \item \textbf{Continuous.}
    The original ODE~\eqref{eq:ssm-continuous}, solved numerically.
\end{enumerate}

\subsection{Diagonalization (S4D)}

The S4D insight is that diagonalizing~$A$ decouples the state into $N$
independent one-dimensional complex systems:
\begin{equation}\label{eq:s4d}
  z_n[t+1] = a_n\,z_n[t] + b_n\,x_n[t],
  \qquad a_n = e^{k_n\dt},
\end{equation}
where $k_n$ is the $n$-th eigenvalue of~$A$.

\subsection{HiPPO initialization and leakage}

The HiPPO-LegS framework provides principled eigenvalues for multi-timescale
memory. For $N$ oscillators remembering a window of length~$T$:
\begin{equation}
  k_n = -\bigl(n+\tfrac{1}{2}\bigr)/T + \im\pi\bigl(n+\tfrac{1}{2}\bigr)/T.
\end{equation}
The real part $\lambda_n = -(n+\tfrac12)/T$ determines the memory timescale
$\tau_n = -1/\lambda_n$; the imaginary part couples frequency to decay so
that each oscillator completes roughly one cycle within its memory window.

In a trainable system, $\lambda$ and $\omega$ are independent parameters:
the SSM theory guarantees a smooth optimization landscape because the kernel
$K[n] = e^{k\dt n}\cdot B$ is smooth in~$k$.

%======================================================================
\section{The Phasor SSM}\label{sec:phasor-ssm}
%======================================================================

Comparing \eqref{eq:rf-layer} with \eqref{eq:s4d} reveals that the
multi-channel R\&F layer \emph{is} a diagonal SSM:

\begin{center}
\begin{tabular}{lll}
\toprule
SSM & R\&F & Role \\
\midrule
$A = \operatorname{diag}(k_1,\dots,k_C)$ &
  $k_c = \lambda_c + \im\omega_c$ & Complex eigenvalues \\
$B = W$ & $W\in\R^{C\times C_{\mathrm{in}}}$ & Input projection \\
$C = I$ & Phase extraction & Readout \\
$D = 0$ & --- & No skip connection \\
\bottomrule
\end{tabular}
\end{center}

Two adjustments are needed to integrate the systems: a
discretization strategy matched to the R\&F neuron's phase-based
communication, and an equal-period constraint that enables direct
spike integration.

\subsection{The equal-period constraint}\label{sec:equal-period}

All layers in a Phasor SSM share the same resonant period~$T$.
When a neuron fires, the receiving neuron oscillates at the same rate,
so the spike arrives as a direct current injection---no intermediate
oscillator is needed to decode a foreign frequency.

If $T_{\text{sender}} \neq T_{\text{receiver}}$, an input oscillator
matched to the sender's frequency would be required to translate the
spike into the receiver's temporal frame. The equal-period constraint
eliminates this stage: the spike's timing within the period already
carries the phase information directly.

\subsection{Two discretization strategies}\label{sec:discretization}

The correct discretization depends on the input type.

\paragraph{Zero-order hold (ZOH)---for continuous inputs.}
When the input is a continuous signal sampled at intervals~$\dt$
(\eg, pixel values encoded as complex numbers), the standard ZOH
discretization applies:
\begin{equation}\label{eq:zoh}
  z_c[n+1] = A_c\,z_c[n] + B_c\sum_j W_{cj}\,x_j[n],
  \qquad A_c = e^{k_c\dt},\quad B_c = \frac{A_c - 1}{k_c}.
\end{equation}

\paragraph{Dirac discretization---for phase/spike inputs.}
When the input is a phase $\theta_j$ from an upstream R\&F layer, it
encodes a spike at time $t_s = (\theta_j/2 + \tfrac12)\,T$ within the
period. The remaining propagation time is
$\delta t_j = T\,(\tfrac12 - \theta_j/2)$.

For a Dirac spike $\delta(t-t_s)$ at input channel~$j$ with
weight~$W_{cj}$, the single-oscillator ODE has the exact solution
\begin{equation}\label{eq:dirac-response}
  z_c(T) = W_{cj}\,e^{k_c\,\delta t_j}.
\end{equation}
This is the free response of a damped oscillator to an impulse: the
exponential $e^{k_c\,\delta t}$ captures both decay
($e^{\lambda_c\,\delta t}$) and rotation ($e^{\im\omega_c\,\delta t}$)
over the remaining time.

Summing over input channels and unrolling over timesteps:
\begin{align}
  H_c[n] &= \sum_j W_{cj}\,e^{k_c\,\delta t_{j}[n]},
    \label{eq:dirac-encoding}\\
  z_c[n+1] &= e^{k_c T}\,z_c[n] + H_c[n].
    \label{eq:dirac-recurrence}
\end{align}
The Dirac kernel uses $K_c[n] = e^{k_c\,n\,T}$ with no ZOH gain factor,
because the input is instantaneous.

\begin{remark}
The Dirac encoding~\eqref{eq:dirac-encoding} depends on the output
channel's eigenvalue~$k_c$ (via $e^{k_c\,\delta t}$), unlike ZOH
encoding which is channel-independent. This is the price of sub-period
phase precision.
\end{remark}

The factored form separates lag-dependent and phase-dependent parts:
\begin{equation}
  e^{k_c\,\Delta t_{\text{elapsed}}}
    = \underbrace{e^{k_c\,\ell\,T}}_{\text{causal kernel}}
    \cdot \underbrace{e^{k_c\,\delta t}}_{\text{Dirac encoding}},
\end{equation}
enabling the convolutional view: the Dirac encoding produces the input
signal~$H_c[n]$, which is then convolved with the causal kernel
$K_c[n] = e^{k_c\,n\,T}$ via FFT.

\subsection{Convolutional view}

Both discretizations produce a linear recurrence
$z_c[n+1] = A_c\,z_c[n] + H_c[n]$, which unrolls to a causal
convolution:
\begin{equation}\label{eq:causal-conv}
  z_c[n] = \sum_{j=0}^{n} K_c[n-j]\,H_c[j]
         = (K_c * H_c)[n],
\end{equation}
where
\[
  K_c[n] =
  \begin{cases}
    e^{k_c\,n\,T} & \text{(Dirac inputs)},\\
    e^{k_c\,n\,\dt}\cdot B_c & \text{(ZOH inputs)}.
  \end{cases}
\]
This convolution is computed via FFT for $O(C\,L\log L\cdot B)$ cost
on long sequences, or via lower-triangular Toeplitz matrix multiplication
for short sequences.

The convolution operates on \textbf{complex-valued} signals, not phases.
Phases enter through the Dirac encoding (which produces complex values
from phase inputs) and are extracted only at the output boundary.

\subsection{Equivalence between views}

The three views are exactly equivalent:
\begin{enumerate}[nosep]
  \item \textbf{Continuous $\to$ Discrete (Dirac).} As the spike kernel
    width $\to 0$ and the ODE solver step $\dd t \to 0$, the spiking ODE
    output converges to the Dirac discrete output.
  \item \textbf{Discrete $\to$ Convolutional.} The causal convolution is
    an algebraic rearrangement of the recurrence.
  \item \textbf{Toeplitz $\to$ FFT.} Zero-padded circular convolution
    reproduces the linear convolution.
\end{enumerate}

\subsection{Phase update dynamics}\label{sec:phase-dynamics}

On the unit circle (after amplitude normalization), the phase update
becomes
\begin{equation}\label{eq:phase-update}
  \theta[n+1]
    = \frac{1}{\pi}\arg\!\bigl(
        A\,e^{\im\pi\theta[n]} + B\,e^{\im\pi\phi[n]}
      \bigr),
\end{equation}
where $A = |A|\,e^{\im\pi\alpha}$ with $\alpha = \omega\dt/\pi$.
This is \textbf{nonlinear in phase} due to the $\arg$ operation on a sum
of complex numbers---fundamentally different from the linear state update
of a standard SSM. The nonlinearity arises because phases combine via
vector addition on the unit circle, not scalar addition.

This has a critical consequence: the convolutional view cannot operate
directly on phases. Computation must proceed in complex space (where
everything is linear) and extract phases only at the output boundary.

\subsection{The frequency diversity--coherence tension}\label{sec:tension}

The filter bank interpretation (Section~\ref{sec:rf-multichannel})
suggests each channel should resonate at a different~$\omega_c$ for
optimal input decomposition. Empirically, networks with trainable
per-channel~$\omega$ achieve higher accuracy.

However, the phase vector
$\bm\theta = (\theta_1,\dots,\theta_C) \in [-1,1]^C$ is only meaningful
as a distributed representation when all channels rotate at the same
rate. If channels have different~$\omega_c$, relative phases drift
through time and similarity measures become time-dependent---the
representation is no longer stable.

Physically, this corresponds to the requirement that neurons in a
population be \textbf{synchronized}: all neurons in a traveling wave share
the same frequency, with phase offsets carrying the representational
content.

\subsection{Resolution: multi-compartment frequency decomposition}
\label{sec:phasorstft}

The \textsc{PhasorSTFT} layer resolves this tension via a multi-compartment
neuron model. Each frequency channel~$f$ has two compartments:

\begin{enumerate}[nosep]
  \item \textbf{Signal compartment} (driven by weighted input):
    \(
      \dd z_{\text{sig}}/\dd t
        = k_f\,z_{\text{sig}} + \sum_j W_{fj}\,I_j(t)
    \)
  \item \textbf{Reference compartment} (free-running, analytically known):
    \(
      z_{\text{ref}}(t) = e^{k_f\,t}
    \)
\end{enumerate}

The \textbf{invariant quantity} is the phase difference:
\begin{equation}
  \Delta\theta_f[n]
    = \arg(z_{\text{sig}}[n]) - \arg(z_{\text{ref}}[n]),
\end{equation}
which represents the input's spectral content at frequency~$\omega_f$,
independent of the carrier. This is the same principle as FM demodulation
or lock-in amplification.

To interface with downstream synchronized layers, the invariant phase is
re-encoded at the downstream carrier~$\omega_{\text{out}}$:
\begin{equation}\label{eq:freq-shift}
  z_{\text{out}}[n]
    = z_{\text{sig}}[n]
    \cdot e^{\im(\omega_{\text{out}} - \omega_f)\,n\,\dt}.
\end{equation}
This frequency-shift modulation is differentiable with respect
to~$\omega_f$, allowing the network to learn which input frequencies are
informative. When $\omega_f = \omega_{\text{out}}$, the shift is the
identity.

The resulting architecture separates concerns:
\begin{equation}
  \text{Input}
    \xrightarrow{\text{PhasorSTFT}}
  \text{frequency-diverse}
    \xrightarrow{\eqref{eq:freq-shift}}
  \text{phase-coherent}
    \xrightarrow{\text{PhasorDense}}
  \text{readout}.
\end{equation}

\begin{remark}[Biological interpretation]
This architecture has a direct analog in dendritic computation.
Pyramidal neuron dendrites sustain local oscillations at frequencies
different from the somatic oscillation; these act as resonant filters
whose signals converge at the soma, modulating the output at the somatic
frequency. The reference compartment corresponds to the dendrite's
intrinsic oscillation; the phase deviation is the information
propagated to the soma.
\end{remark}

%======================================================================
\section{Connections to Holographic Distributed Representations}
\label{sec:hd}
%======================================================================

The Phasor SSM's phase vectors are instances of
\textbf{Fourier Holographic Reduced Representations} (Fourier HRR),
the frequency-domain formulation of Plate's HRR algebra.
In a Fourier HRR, symbols are represented as complex vectors on the unit
torus $\bm{s} = (e^{\im\pi\theta_1},\dots,e^{\im\pi\theta_C}) \in \C^C$
with $|\bm{s}_c| = 1$. Two algebraic operations define the representational
system.

\subsection{Binding as phase addition (complex multiplication)}

The VSA \textbf{binding} operation corresponds to element-wise
complex multiplication:
\begin{equation}\label{eq:bind}
  \texttt{bind}(\bm\theta_1, \bm\theta_2)_c
    = \operatorname{remap}\bigl(\theta_{1,c} + \theta_{2,c}\bigr),
  \qquad
  e^{\im\pi\theta_{1,c}} \cdot e^{\im\pi\theta_{2,c}}
    = e^{\im\pi(\theta_{1,c}+\theta_{2,c})}.
\end{equation}
Binding creates a new representation that is dissimilar to both operands
but can be inverted (by binding with the conjugate), enabling
variable--value association.

In the R\&F context, the per-step state rotation
$A_c = e^{k_c\dt}$ is itself a binding with the natural phase increment
$\alpha_c = \omega_c\dt/\pi$. The SSM's temporal dynamics therefore
implement a \textbf{sequence of bindings} with a fixed time-phase:
\begin{equation}
  z_c[n] = \texttt{bind}(z_c[0],\, n\cdot\alpha_c)
            + \text{input contributions}.
\end{equation}
The accumulated rotation $n\cdot\alpha_c$ is the natural oscillator
phase at step~$n$, and input contributions are phase-shifted copies
of the input, each rotated by the appropriate number of time steps.

\subsection{Bundling as superposition (complex addition)}

The VSA \textbf{bundling} operation corresponds to element-wise complex
addition followed by phase extraction:
\begin{equation}\label{eq:bundle}
  \texttt{bundle}(\bm\theta_1, \bm\theta_2)_c
    = \frac{1}{\pi}\arg\!\bigl(
        e^{\im\pi\theta_{1,c}} + e^{\im\pi\theta_{2,c}}
      \bigr).
\end{equation}
Bundling is the constructive/destructive interference of unit phasors:
similar phases reinforce, dissimilar phases cancel. This gives the
R\&F system an inherent \textbf{similarity-weighted averaging} that
linear SSMs lack.

In the SSM recurrence, bundling occurs at every step:
\begin{equation}
  z[n+1] = A\,z[n] + B\,I[n]
  \qquad\text{(complex addition = bundling of rotated state and new input)}.
\end{equation}
The causal convolution unrolls this as a \textbf{temporal bundling}:
\begin{equation}\label{eq:temporal-bundle}
  z_c[n] = \sum_{j=0}^{n} K_c[n-j]\,H_c[j]
  \qquad\text{(superposition of all past inputs at different decay/rotation)}.
\end{equation}
Each past input is bound with a time-dependent phase rotation
(via the kernel~$K_c$) and weighted by exponential decay, then all
contributions are bundled into the current state.

\subsection{The R\&F recurrence as VSA composition}

Combining binding and bundling, the phase update~\eqref{eq:phase-update}
decomposes as:
\begin{equation}\label{eq:vsa-recurrence}
  \theta[n+1]
    = \texttt{bundle}\!\bigl(
        \texttt{bind}(\theta[n],\,\alpha),\;
        \texttt{bind}(\phi[n],\,\beta)
      \bigr),
\end{equation}
where $\alpha$ is the per-step rotation from~$A$ and
$\beta = \arg(B)/\pi$ is the input phase shift from the discretization
gain.

\begin{proposition}
The R\&F recurrence is a composition of Fourier HRR binding
(phase rotation) and bundling (interference), not a linear recurrence
in phase space. The linearity of the underlying complex-space dynamics
is what enables the convolutional view; the nonlinearity enters only at
the phase-extraction boundary.
\end{proposition}

\subsection{Temporal holography: sequence memory via convolution}

The causal convolution~\eqref{eq:temporal-bundle} has a direct
interpretation as a \textbf{holographic sequence memory}.
In the HRR literature, a sequence $x_0, x_1, \dots, x_{L-1}$
is stored as the superposition
\begin{equation}\label{eq:hrr-trace}
  S = \sum_{n=0}^{L-1} \texttt{bind}(x_n,\, p^n),
\end{equation}
where $p$ is a position/time vector and $p^n$ denotes $n$-fold
self-binding.

In the Phasor SSM, the per-step rotation $A_c = e^{k_c T}$ plays the
role of the position code~$p$ (with $p^n = A_c^n = e^{k_c\,nT}$),
and the convolution kernel $K_c[n] = A_c^n$ implements the
time-binding automatically. The convolution output is therefore
a holographic trace: each input is bound with its temporal position
and all are bundled into the state vector.

The exponential decay $|A_c^n| = e^{\lambda_c\,nT}$ provides a
natural \textbf{recency weighting}: recent inputs dominate the
hologram, older inputs fade. The decay rate $\lambda_c$ controls
the memory horizon, directly linking the SSM's leakage parameter
to the HRR trace's temporal resolution.

\subsection{Similarity-based readout as holographic retrieval}

The readout equation in the Phasor SSM computes cosine similarity
between the state phase vector and a codebook of prototypes:
\begin{equation}\label{eq:similarity-readout}
  y_k = \frac{1}{C}\sum_{c=1}^{C}
    \cos\!\bigl(\pi(\theta_c - \theta_c^{(k)})\bigr),
\end{equation}
where $\bm\theta^{(k)}$ is the $k$-th codebook entry.
This is exactly the \textbf{Fourier HRR similarity} measure:
the real part of the normalized inner product between two unit-modulus
complex vectors,
\[
  y_k = \frac{1}{C}\operatorname{Re}\!\Bigl(
    \sum_c e^{\im\pi\theta_c}\,\overline{e^{\im\pi\theta_c^{(k)}}}
  \Bigr)
  = \frac{1}{C}\sum_c \cos\bigl(\pi(\theta_c - \theta_c^{(k)})\bigr).
\]
Classification therefore reduces to holographic retrieval:
the network's temporal bundling produces a query vector, and the
codebook lookup finds the nearest prototype in the HRR space.

\subsection{The equal-period constraint as algebraic consistency}

The equal-period constraint (Section~\ref{sec:equal-period}) has a
natural interpretation in the HRR framework. For the phase vector
$\bm\theta\in[-1,1]^C$ to be a valid holographic representation,
all components must share the same temporal frame---otherwise the
binding/unbinding algebra is inconsistent.

If channel~$c$ rotates at~$\omega_c$, then after $n$~steps the
time code is $e^{\im\omega_c\,n\,T}$, which differs per channel.
Unbinding a temporal position (to retrieve what was stored at time~$j$)
would require channel-dependent position codes, breaking the
element-wise structure of the HRR.

When all channels share $\omega$, the time code $p = e^{\im\omega T}$
is a scalar applied uniformly, and the HRR algebra is well-defined:
binding, unbinding, and similarity all operate element-wise with a
consistent temporal reference.

The \textsc{PhasorSTFT} layer (Section~\ref{sec:phasorstft}) resolves
the tension between frequency diversity and algebraic consistency by
confining per-channel~$\omega$ to the input decomposition stage and
re-encoding at a uniform carrier before the representation enters the
HRR regime.

%======================================================================
\section{Formal System Summary}\label{sec:summary}
%======================================================================

\begin{definition}[Phasor SSM]
A Phasor SSM is a diagonal complex state-space model with parameters
$W\in\R^{C\times C_{\mathrm{in}}}$,
$\lambda\in\R^C_-$, $\omega\in\R^C_+$, $T\in\R_+$,
and derived quantities
$k_c = \lambda_c + \im\omega_c$,
$A_c = e^{k_c T}$,
$B_c = (A_c-1)/k_c$.
It admits three equivalent computational views:

\medskip
\noindent\textbf{Continuous} (direct spike integration):
\[
  \frac{\dd z_c}{\dd t} = k_c\,z_c(t) + \sum_j W_{cj}\,I_j(t),
  \qquad I_j(t) = \sum_s \kappa(t-t_s).
\]

\noindent\textbf{Discrete} (Dirac for phase inputs, ZOH for continuous):
\[
  z_c[n+1] = A_c\,z_c[n] + H_c[n],
  \qquad H_c[n] =
  \begin{cases}
    \sum_j W_{cj}\,e^{k_c\,\delta t_j[n]} & \text{(Dirac)},\\
    B_c\sum_j W_{cj}\,x_j[n] & \text{(ZOH)}.
  \end{cases}
\]

\noindent\textbf{Convolutional}:
\[
  Z_c = K_c * H_c, \qquad
  K_c[n] = A_c^n\ (\text{Dirac}) \text{ or } A_c^n B_c\ (\text{ZOH}).
\]

\medskip
\noindent The phase readout $\theta_c = \arg(z_c)/\pi$ introduces a
nonlinearity equivalent to Fourier HRR bundling. The temporal dynamics
implement HRR binding with a time-phase. The codebook readout implements
HRR similarity retrieval.
\end{definition}

\bigskip

\begin{center}
\renewcommand{\arraystretch}{1.2}
\begin{tabular}{@{}lll@{}}
\toprule
\textbf{Concept} & \textbf{Standard SSM} & \textbf{Phasor SSM} \\
\midrule
State variable & $u\in\R^N$ or $\C^N$ & $z\in\C^C$ (complex potential) \\
State transition & $A\in\R^{N\times N}$ & $\operatorname{diag}(e^{k_c T})$
  (diagonal) \\
Input projection & $B\in\R^{N\times d}$ &
  ZOH: $\operatorname{diag}(B_c)\cdot W$;\;
  Dirac: $e^{k_c\delta t}\cdot W$ \\
Readout & $C\,u$ (linear) & $\arg(\cdot)/\pi$ (nonlinear phase) \\
State addition & Linear sum & Bundling (interference) \\
Time evolution & $e^{A\dt}$ & Binding (phase rotation) \\
Similarity & Inner product & $\cos(\pi(\theta-\theta_{\text{ref}}))$
  (Fourier HRR) \\
Sequence memory & Recurrence & Holographic trace~\eqref{eq:hrr-trace} \\
Frequency analysis & Fixed & PhasorSTFT (trainable $\omega$, re-encoded) \\
\bottomrule
\end{tabular}
\end{center}

\end{document}
